% Vietnamese translation for OI Public Library
% Source-PDF: 英文资料 ENGLISH MATERIALS/A Probability Path - Resnick.pdf
\documentclass[11pt]{article}
\usepackage{oipl-vn}

\title{Một con đường vào xác suất}
\author{Sidney I. Resnick}
\date{Birkhauser, 1999}

\begin{document}
\maketitle

\origtitle{A Probability Path}
\sourcepath{English Materials / A Probability Path - Resnick.pdf}

\renewcommand{\abstractname}{Tóm tắt}
\begin{abstract}
Đây là ghi chú dịch tiếng Việt cho sách \emph{A Probability Path} của Sidney
I. Resnick. Bản PDF có lớp OCR đọc được và mục lục đầy đủ. Ghi chú này chuyển
ngữ phần thông tin đầu sách, mục lục chương và vai trò của từng chương, nhằm
giúp người học thuật toán xác định các phần xác suất--độ đo cần đọc sâu hơn.
\end{abstract}

\section{Thông tin sách}

\begin{center}
\begin{tabular}{ll}
\textbf{Tác giả} & Sidney I. Resnick \\
\textbf{Cơ quan} & Cornell University \\
\textbf{Nhà xuất bản} & Birkhauser Boston \\
\textbf{Năm} & 1999 \\
\textbf{ISBN} & 0-8176-4055-X
\end{tabular}
\end{center}

Sách trình bày xác suất hiện đại theo hướng dựa trên độ đo. Nội dung phù hợp
với người đã quen giải tích cơ bản và muốn xây dựng nền tảng chắc chắn cho
biến ngẫu nhiên, kỳ vọng, hội tụ, định lý giới hạn và martingale.

\section{Mục lục chương đã dịch}

\begin{description}
  \item[Chương 1. Tập hợp và biến cố]
  Nhắc lại lý thuyết tập hợp, giới hạn của dãy tập, các lớp tập đóng dưới phép
  toán, trường sigma sinh bởi một lớp tập và tập Borel trên đường thẳng thực.

  \item[Chương 2. Không gian xác suất]
  Định nghĩa và tính chất cơ bản của không gian xác suất, các định lý đóng,
  định lý Dynkin, xây dựng độ đo xác suất, độ đo Lebesgue trên \((0,1]\) và
  xây dựng xác suất từ hàm phân phối.

  \item[Chương 3. Biến ngẫu nhiên, phần tử ngẫu nhiên và ánh xạ đo được]
  Trình bày ảnh ngược, ánh xạ đo được, phân phối cảm sinh, phần tử ngẫu nhiên
  trong không gian metric, tính đo được qua giới hạn và trường sigma sinh bởi
  ánh xạ.

  \item[Chương 4. Độc lập]
  Định nghĩa độc lập cho biến cố và biến ngẫu nhiên, các ví dụ về record và
  biểu diễn nhị phân, nhóm các biến độc lập, bổ đề Borel--Cantelli và các định
  luật không--một.

  \item[Chương 5. Tích phân và kỳ vọng]
  Xây dựng tích phân qua hàm đơn giản, định nghĩa kỳ vọng, các tính chất cơ
  bản, định lý giới hạn cho tích phân, mật độ, so sánh tích phân Riemann và
  Lebesgue, không gian tích và định lý Fubini.

  \item[Chương 6. Các khái niệm hội tụ]
  Hội tụ hầu chắc chắn, hội tụ theo xác suất, hội tụ trong \(L^p\), tích phân
  đều, các bất đẳng thức chuẩn và quan hệ giữa các dạng hội tụ.

  \item[Chương 7. Luật số lớn và tổng các biến ngẫu nhiên độc lập]
  Kỹ thuật cắt cụt, luật yếu tổng quát, hội tụ hầu chắc chắn của tổng độc lập,
  luật mạnh số lớn và định lý ba chuỗi Kolmogorov.

  \item[Chương 8. Hội tụ theo phân phối]
  Định nghĩa hội tụ yếu, bổ đề Scheffe, định lý Skorohod dạng cơ bản, phương
  pháp delta, định lý Portmanteau, các quan hệ giữa dạng hội tụ và ứng dụng
  cho phân phối cực trị.

  \item[Chương 9. Hàm đặc trưng và định lý giới hạn trung tâm]
  Ôn hàm sinh moment, định nghĩa hàm đặc trưng, khai triển, moment và đạo hàm,
  định lý duy nhất, định lý liên tục, độ chặt, định lý Prohorov, CLT cổ điển
  và CLT Lindeberg--Feller.

  \item[Chương 10. Martingale]
  Định lý Radon--Nikodym, kỳ vọng có điều kiện, martingale, submartingale,
  thời điểm dừng, upcrossing, các định lý hội tụ, đồng nhất thức Wald, bước
  đi ngẫu nhiên, martingale ngược và một phần nhập môn tài chính toán học.
\end{description}

\section{Ghi chú nội dung}

Sách không phải là tài liệu OI trực tiếp, nhưng nhiều chương cung cấp nền tảng
cho các chủ đề thường gặp trong phân tích thuật toán ngẫu nhiên:
\begin{itemize}
  \item Chương 1--3 làm rõ ngôn ngữ độ đo, biến ngẫu nhiên và phân phối.
  \item Chương 4--7 là lõi cho xác suất rời rạc nâng cao, biến độc lập, kỳ vọng
        và luật số lớn.
  \item Chương 8--9 hỗ trợ đọc các chứng minh hội tụ, xấp xỉ phân phối và định
        lý giới hạn trung tâm.
  \item Chương 10 hữu ích cho các bài toán quá trình ngẫu nhiên, thời điểm dừng
        và martingale.
\end{itemize}

\section{Ghi chú về trích xuất}

Tệp PDF có 459 trang. Lớp OCR trích xuất được văn bản tiếng Anh tương đối tốt,
nhưng có lỗi nhận dạng nhỏ ở ký hiệu toán và một vài từ trong trang bìa. Bản
dịch đầy đủ từng chương cần đối chiếu PDF gốc khi chuyển công thức, chỉ số và
ký hiệu xác suất sang LaTeX.

\end{document}
